\documentclass[11pt,a4paper]{article}
\usepackage[margin=1in]{geometry}
\usepackage[T1]{fontenc}
\usepackage{lmodern}
\usepackage{microtype}
\usepackage{amsmath,amssymb}
\usepackage{xcolor}
\usepackage{hyperref}
\hypersetup{colorlinks=true,linkcolor=blue!50!black,urlcolor=blue!50!black}
\title{Compiling the Raytracer to the Native CPU}
\author{Generated C, x86-64 assembly, executable, and PPM output}
\date{16 August 2026}
\begin{document}
\maketitle
\begin{abstract}
This note adds a native-host backend to the parametric raytracer compiler. A raytracer semantic description is quoted and unquoted, specialized directly into x86-64 AT\&T assembly, assembled with `as --64`, linked with `ld`, and run. The executable produces and validates a 96 by 64 binary PGM image. No C compiler is used for the target build. This backend is intentionally separate from the AArch64 emitter: the high-level ray semantics are shared while the final ABI and machine lowering are target-specific.
\end{abstract}
\section{Compilation pipeline}
The tested pipeline is
\[
D_{ray}\xrightarrow{\mathsf{quote}/\mathsf{unquote}}x86\text{-}64\text{ assembly}
\xrightarrow{\texttt{as --64}}\text{object}
\xrightarrow{\texttt{ld}}\text{executable}
\xrightarrow{}\texttt{render.ppm}.
\]
The generated C contains a fixed-point orthographic sphere test and a nested pixel loop. It writes an ASCII PPM image with a standard header and one grayscale pixel for each coordinate.

\section{Semantic specialization}
The description contains width, height, sphere center, radius, and output path. The generator emits these as static constants in the C source. Direct generation and generation through a quoted description are compared byte-for-byte before invoking the host compiler. Therefore the tested path is genuinely staged rather than merely two independently written generators.

\section{Native build and execution}
The implementation is \texttt{cpu\_ray\_codegen.cpp}. The exact commands used were:
\begin{verbatim}
g++ -std=c++23 -Wall -Wextra -pedantic -O2 cpu_ray_codegen.cpp -o cpu_ray_codegen
./cpu_ray_codegen > generated_raytracer.s
as --64 generated_raytracer.s -o generated_raytracer.o
ld generated_raytracer.o -o generated_raytracer
./generated_raytracer
\end{verbatim}
The resulting executable is a statically linked x86-64 ELF binary on the current machine. Running it succeeds and produces a valid PGM file whose header reports $96\times64$ pixels and maximum value 255. The output file is non-empty and was identified as Netpbm image data.

\section{Relation to the AArch64 backend}
The AArch64 backend emits a ray kernel and an explicit pixel sink ABI. The native backend chooses generated C as its portable intermediate representation and delegates final register allocation, calling convention, and object emission to the host C compiler. This is a useful staged bootstrap path: the same semantic description can target both a direct AArch64 emitter and a native CPU toolchain without conflating their ABIs.

\section{Limitations}
The current renderer uses a binary PGM sink and a single sphere with fixed-point arithmetic. It does not yet implement floating-point vectors, multiple objects, materials, shadows, or acceleration structures. Those extensions can be added to the semantic description while preserving the quote/unquote boundary and the output-image validation step.
\end{document}
